% Clinical Background section. Paste into your own thesis, or \input it.

\section{Coronary computed tomography angiography}
\label{sec:ccta}

Coronary computed tomography angiography (CCTA) is the non-invasive imaging test recommended
as an initial investigation in patients with suspected chronic coronary syndrome
\cite{knuuti2019esc}.
An iodinated contrast agent is injected and the scan is timed so that acquisition coincides
with peak opacification of the coronary lumen.
The result is a three-dimensional volume in which the contrast-filled lumen is bright against
surrounding tissue.

This section treats CCTA as a measuring instrument, because the limits of the instrument
bound any geometric claim made from these images.

\subsection{Acquisition and its failure modes}

The coronary arteries move throughout the cardiac cycle, so acquisition is synchronized to
the electrocardiogram and data are reconstructed at the quiescent phase, conventionally in
mid-diastole \cite{abbara2016scct}.
Heart rate is commonly lowered pharmacologically before the scan to lengthen that quiescent
window \cite{abbara2016scct}.
When synchronization fails, or when the rhythm is irregular, motion artifact blurs or
duplicates the vessel \cite{abbara2016scct}.

Dense calcium produces blooming, in which the apparent extent of a calcified plaque exceeds
its true extent and the adjacent lumen is obscured \cite{abbara2016scct}.
Blooming is the reason a heavily calcified segment may be uninterpretable even in an
otherwise excellent scan.

\subsection{Resolution against vessel caliber}

The ImageCAS-X volumes have in-plane spacing between 0.318 mm and 0.43 mm and a
slice thickness of 0.5 mm, varying between cases \cite{zeng2023imagecas}.
Proximal coronary arteries are a few millimeters in diameter and taper distally toward 1 mm
and below.
A vessel two or three voxels across cannot support a confident radius measurement, and this
places a floor on the branch order at which geometric features remain meaningful.

Contrast opacification in this cohort is adequate: sampling the volumes at the annotated
centerline points, after converting them from the centerlines' LPS frame into voxel indices,
gives a median of 336 HU, consistent with a contrast-filled lumen.

\subsection{Image quality as a variable}

Image quality is graded from 0 to 4 in this dataset, and the 200 cases graded 0 are excluded
from the released cohort, leaving the 800 used here \cite{bransby2026imagecasx}.
Among those 800 the grading remains informative rather than uniform.
Because a lower-quality scan plausibly yields fewer traceable distal branches, image quality
is a candidate confounder for any branch-count statistic derived from this cohort.
